\section*{Capítulo \ref{ch:g9:landforms} --- As formas de relevo}

\begin{solution}{exo:g9:landforms:1}
(a) O vento. (b) A água que dissolve a \omterm{def:g3:stones-pebbles-sand:rock}{rocha} (e deposita calcita na
\omterm{def:g9:landforms:karst}{caverna}). (c) O \omterm{def:g2:water-on-earth:ocean}{mar} (\omterm{def:g9:landforms:beach}{deriva litorânea}). (d) O gelo. (e) O \omterm{def:g2:water-on-earth:ocean}{mar}. (f) Um \omterm{def:g1:landscapes-around-us:water}{rio}.
\end{solution}

\begin{solution}{exo:g9:landforms:2}
(a) \omterm{def:g9:landforms:uplift}{Soerguimento}. (b) A \omterm{def:g9:landforms:dune}{face de avalanche}. (c) Uma \omterm{def:g9:landforms:karst}{dolina}. (d) A \omterm{def:g9:landforms:beach}{deriva
litorânea}. (e) Um \omterm{def:g9:landforms:karst}{carste}.
\end{solution}

\begin{solution}{exo:g9:landforms:3}
(a) Falso: os chifres apontam no sentido do vento; o nariz arredondado
fica de frente para o vento. (b) Verdadeiro. (c) Falso: o vento a move
grão por grão, encosta suave acima e \omterm{def:g9:landforms:dune}{face de avalanche} abaixo. (d)
Verdadeiro.
\end{solution}

\begin{solution}{exo:g9:landforms:4}
(b), (d), (e), (a), (c).
\end{solution}

\begin{solution}{exo:g9:landforms:5}
(a) Do oeste: os chifres apontam no sentido do vento. (b) No lado leste,
entre os chifres. (c) $60 \div 6 = 10$ anos. (d) Uma \omterm{def:g9:landforms:dune}{duna} maior contém
mais \omterm{def:g3:stones-pebbles-sand:sand}{areia}, e toda ela tem de passar por cima da crista para a \omterm{def:g9:landforms:dune}{duna}
avançar.
\end{solution}

\begin{solution}{exo:g9:landforms:6}
(a) $7 \times 1 = \qty{7}{m}$. (b) $2 \div 12 = \frac{1}{6}$ de um ano.
(c) Sete anos de recuo em $\frac{1}{6}$ de um ano: $7 \times 6 = 42$
vezes mais depressa. (d) As ondas primeiro escavam um entalhe; a \omterm{def:g3:stones-pebbles-sand:rock}{rocha} de
cima fica em balanço sobre ele por algum tempo, sem se mexer; depois a
parte em balanço desaba de uma só vez, e a borda salta para trás.
\end{solution}

\begin{solution}{exo:g9:landforms:7}
(a) Para leste. (b) Três trajetos \omterm{def:g9:landforms:beach}{praia} acima, inclinados para o
nordeste, cada um seguido de um trajeto direto de volta para o sul: um
zigue-zague que avança para leste. (c) Um \omterm{def:g9:landforms:beach}{pontal}, crescendo para leste
através da entrada da baía. (d) O \omterm{def:g9:landforms:beach}{pontal} abriga das ondas a água que fica
atrás dele; em água calma, a \omterm{def:g3:stones-pebbles-sand:sand}{lama} fina pode se depositar.
\end{solution}

\begin{solution}{exo:g9:landforms:8}
(a) Ela absorve gás carbônico do ar e do solo. (b) O ácido fraco dissolve
o calcário devagar, mas quase não ataca o granito. (c) A água só pode
entrar na \omterm{def:g3:stones-pebbles-sand:rock}{rocha} pelas suas fendas e ao longo das suas camadas; ela alarga
primeiro essas aberturas. (d) A \omterm{def:g3:stones-pebbles-sand:rock}{rocha} embaixo da superfície se dissolve e
o terreno afunda devagar; ou o teto de uma \omterm{def:g9:landforms:karst}{caverna} desaba.
\end{solution}

\begin{solution}{exo:g9:landforms:9}
(a) Quando a gota encontra o ar da \omterm{def:g9:landforms:karst}{caverna}, ela libera parte do seu gás
carbônico e não consegue mais manter todo o calcário dissolvido, que sai
na forma de calcita. (b) Uma estalagmite cresce a partir de gotas que
respingam no chão e se espalham; não há gota pendurada na sua ponta para
construir um tubo.
(c) $\qty{30}{cm} = \qty{300}{mm}$; $300 \div 0.1 = 3000 = 3 \times 10^3$
anos.
\end{solution}

\begin{solution}{exo:g9:landforms:10}
(a) $4600 \div 10^7 = 4.6 \times 10^{-4}$~m por ano, ou \qty{0.46}{mm}
por ano. (b) $0.46 - 0.2 = \qty{0.26}{mm}$ por ano. (c) $\qty{1}{m} =
\qty{1000}{mm}$; $1000 \div 0.46 \approx 2200$, cerca de duas mil vezes
mais rápido. (d) Ele é mantido por milhões de anos: $\qty{0.26}{mm} \times
10^7 = 2.6 \times 10^6$~mm, cerca de \qty{2.6}{km}.
\end{solution}

\begin{solution}{exo:g9:landforms:11}
(a) Forma: \omterm{def:g9:landforms:dune}{dunas} em crescente com a \omterm{def:g9:landforms:dune}{face de avalanche} no lado sudeste.
Agente e processo: o vento, soprando do noroeste, \omterm{def:g5:water-shapes-the-land:deposition}{deposição}. Material:
\omterm{def:g3:stones-pebbles-sand:sand}{areia} solta. (b) Forma: uma \omterm{def:g9:landforms:cliff}{falésia} acima de uma laje plana descoberta na
maré baixa. Agente e processo: o \omterm{def:g2:water-on-earth:ocean}{mar}, \omterm{def:g5:water-shapes-the-land:erosion}{erosão}; a laje é uma plataforma de
abrasão deixada pela \omterm{def:g9:landforms:cliff}{falésia} que recua. Material: \omterm{def:g3:stones-pebbles-sand:rock}{rocha} branca,
provavelmente giz, uma \omterm{def:g3:stones-pebbles-sand:rock}{rocha} mole. (c) Forma: um riacho que some num
buraco, uma \omterm{def:g5:water-shapes-the-land:source-mouth}{nascente} longe dali. Agente e processo: água que dissolve a
\omterm{def:g3:stones-pebbles-sand:rock}{rocha} e corre embaixo da terra. Material: calcário; o planalto é um
\omterm{def:g9:landforms:karst}{carste}.
\end{solution}

\begin{solution}{exo:g9:landforms:12}
(a) A \omterm{def:g9:landforms:dune}{duna} pequena: $45 \div 6 = 7.5$ anos; a maior: $45 \div 3 =
15$ anos. (b) $45 \div 100 = 0.45$ de um ano, cerca de cinco ou seis
meses. (c) Por exemplo, plantar capim nas \omterm{def:g9:landforms:dune}{dunas}: as suas raízes seguram a
\omterm{def:g3:stones-pebbles-sand:sand}{areia}, e o vento não consegue mais levá-la.
\end{solution}

\begin{solution}{pb:g9:landforms:1}
\textbf{1.} $46.0 - 30.4 = \qty{15.6}{m}$.

\textbf{2.} $15.6 \div 20 = \qty{0.78}{m}$ por ano.

\textbf{3.} Entre 2012 e 2014: $41.0 - 36.2 = \qty{4.8}{m}$ em dois
anos, ou $4.8 \div 2 = \qty{2.4}{m}$ por ano.

\textbf{4.} $2.4 \div 0.78 \approx 3$ vezes mais rápido. Provavelmente
um inverno de \omterm{def:g5:natural-hazards:storm}{tempestades}, durante o qual grandes partes em balanço
desabaram.

\textbf{5.} \qty{78}{cm} por ano; $0.78 \times 100 = \qty{78}{m}$ por
século.

\textbf{6.} Sim: \qty{0.78}{m} por ano é um pouco menos que um metro por ano.

\textbf{7.} $17 \div 0.78 \approx 22$ anos.

\textbf{8.} A \omterm{def:g9:landforms:cliff}{falésia} recua aos saltos: um intervalo de dois anos pode
conter um grande desmoronamento, ou nenhum. Ao longo de 20 anos, os
saltos e os anos calmos se compensam.

\textbf{9.} A borda ainda tem $30.4 - 10 = \qty{20.4}{m}$ a percorrer: $20.4 \div 0.78
\approx 26$ anos, ou seja, em 2050.

\textbf{10.} $2 \times 0.78 = \qty{1.56}{m}$ por ano; $30.4 \div 1.56
\approx 19.5$ anos depois de 2024, ou seja, em 2043.

\textbf{11.} A borda de uma \omterm{def:g9:landforms:cliff}{falésia} pode ceder a qualquer momento; ficar
em cima dela, ou perto dela, é perigoso.

\textbf{12.} $30.4 \div 0.78 \approx 39$ anos depois de 2024: em
\textbf{2063}.
\end{solution}
